\section{Background and Related Literature}\label{sec:background}

The fiscal weight of Scotland's labour market rests on the settlement that followed the Scotland Acts of 2012 and 2016. Devolved and partially devolved taxes now fund close to half of the Scottish Government's resource spending, with non-savings, non-dividend income tax the single largest component \citep{sfc2026}. The block grant adjustment (BGA) is what makes the composition of the tax base consequential. When income tax was devolved the block grant was cut by the revenue foregone, and that deduction indexed thereafter to the growth of the equivalent rUK tax base \citep{gov2023}. The Scottish budget therefore gains only insofar as Scottish income tax grows faster than its rUK comparator, so that the level of Scottish revenue is immaterial and only the differential matters. The Scottish Fiscal Commission (SFC) records this as the income tax net position, the gap between devolved revenue and the BGA, forecast at £969 million for 2026--27 \citep{sfc2026}. A relative deterioration in Scottish earnings, however modest in absolute terms, feeds mechanically into a smaller net position, and because outturn is reconciled against forecast three years on, an error made today constrains spending well into the future \citep{gov2023}.

It is against this backdrop that the fiscal repercussions of AI-driven labour-market disruption matter. A technology that raises productivity and reshapes labour demand unevenly across occupations will reshape earnings unevenly across the economy, and any resulting divergence between the Scottish and rUK tax bases surfaces in the net position the SFC must forecast \citep{sfc2016}. There is good reason to expect such a divergence, rooted in Scotland's distinctive occupational composition, and its direction is an empirical question: the composition may harbour Scotland from displacement, or it may deny Scotland the gains that accrue to complemented work. This is heightened by a longstanding productivity gap. With growth weak across Scotland, the UK, and much of Europe since the 2008 crisis, and forecast downgrades to productivity already worsening the fiscal outlook \citep{sfc2026}, whether AI narrows or widens that gap turns on how its gains are distributed across regions.

\subsection*{The task mechanism: displacement, complementarity, and the cognitive frontier}

The labour-market consequences of AI run through the task. Whether AI raises or lowers the demand for a worker depends on which of the activities that make up the job the technology can complete, complement, or cannot do. The organising framework is the task-based model of \citet{acemoglu2019}, in which production allocates a continuum of tasks between labour and capital, and a technology acts through two opposing channels. A displacement effect operates as capital absorbs tasks workers performed, and a reinstatement effect operates as new tasks restore a comparative advantage for labour. The central lesson is that productivity gains and labour demand can move apart, and \citet{acemoglu2020} show that ``so-so'' automation, which displaces workers while delivering only modest productivity gains, can depress employment and the labour share even as output rises. Exposure to a technology is thus not the same as harm from it: a task AI can perform is displaced, whereas a task AI makes a worker faster at is complemented, and the two carry opposite signs for labour demand.

Which tasks a technology reaches has changed. Earlier computerisation acted on routine, codifiable tasks, the empirical content of the routine-biased technical change literature \citep{autor2003, autor2013}, hollowing out middle-wage clerical and production work while sparing both manual and analytical non-routine activity. Generative AI reaches the non-routine cognitive tasks the previous wave could not, so exposure no longer maps onto the routine/non-routine divide that organised two decades of work. The theoretical fault line is consequently the substitution--complementarity margin, and the literature divides on which dominates. \citet{autor2024} reads AI optimistically, as a tool that could reinstate middle-skill judgement and widen access to expert tasks, whereas \citet{acemoglu2025} is more cautious, projecting an aggregate total-factor-productivity gain of under one per cent over a decade once only currently automatable tasks are counted. The disagreement concerns magnitude and net sign, not the mechanism, and it remains unresolved.

\subsection*{Measuring exposure and the reliability problem}

Turning the mechanism into a measure requires judging, task by task, what the technology can do, at a scale expert assessment cannot reach. Successive generations of measures bet on different inputs: engineering bottlenecks read at the level of the whole occupation \citep{frey2017}, an approach later criticised for treating occupations as indivisible; the abilities an occupation requires \citep{felten2021}; the textual overlap of patents with job descriptions \citep{webb2019}; and the task itself, scored for its suitability for machine learning \citep{brynjolfsson2018}. Large language models then changed the economics of the measurement. \citet{eloundou2024} have an LLM rate each O*NET task against a capability rubric, and show the ratings track expert annotators closely while covering the whole taxonomy, which no human exercise can afford. The approach now anchors institutional work: \citet{henseke2025} build a task-based UK index from LLM ratings, the \citet{obr2025} score every O*NET task to feed a national productivity projection, and the IMF's complementarity-adjusted index separates displacement from augmentation \citep{pizzinelli2023, cazzaniga2024}.

These measures are not interchangeable. Since they read exposure off different inputs (bottlenecks, abilities, patents, or task text) they rank the same occupation differently, and any estimate built on one inherits that disagreement. The reliability problem has lately sharpened from a cross-measure concern into a within-measure one. \citet{yin2026} score an identical task set with four contemporary language models and find average exposure differs by more than a factor of three across them, with downstream labour-market magnitudes varying roughly 2.4-fold by model. The divergence is systematic as each model carries a persistent tendency to score high or low, such that averaging does not remove it. That exposure is an indicative quantity, and not the same thing as realised impact, is now firm across the institutional literature. What remains under-developed is any propagation of that measurement uncertainty into the estimates built on top of it.

\subsection*{The geography of exposure and Scotland's composition}

Exposure is a property of tasks and tasks bundle into occupations, so AI's incidence is uneven across space, though the literature disagrees on how uneven at the spatial scales of policy interest. \citet{tbi2024} find limited regional variation in UK exposure, attributing it to the diversified sectoral mix of most regions, where the OECD finds cross-country exposure ranging from 16--77\% with complementarity rising in regional education \citep{oecd2024}, and European evidence finds automation risk driven systematically by industrial structure and agglomeration \citep{crowley2021}. The disagreement is partly artefactual: coarse administrative geographies average over the functional labour markets where local conditions actually operate \citep{overman2022}, so the regional signal shows up most clearly in occupational mix and least clearly in coarse exposure averages. Consistent with this, \citet{henseke2025} attribute the recent rise in UK exposure to shifts in occupational composition rather than to within-occupation task change.

Scotland's occupational mix is distinctive in ways that bear directly on exposure. Its public sector accounts for around 22\% of employment against roughly 18\% UK-wide, a gap persistent over time and larger than that of any English region, and services account for 77.1\% of Scottish GDP against 80.5\% for the UK, with more local-government delivery of services such as social care and correspondingly less weight in the professional, scientific and information services concentrated in London and the South East \citep{scotgov2024, commonslibrary2024}. This composition is double-edged. Financial, insurance and professional services are among the most exposed parts of the UK economy, and caring, elementary and skilled-trade work among the least \citep{obr2025}. But the exposed sectors divide on the substitution--complementarity margin, since public administration, health and caring work scores high on complementarity and low on substitutability, whereas professional and financial work is exposed on both. Whether Scotland's lighter concentration in the latter harbours its workforce or merely denies it the productivity gains of complemented work is therefore ambiguous, and resolvable only empirically. The fiscal edge is sharper still because public-sector pay in Scotland sits about 5\% above the UK median and has diverged upward since 2019 \citep{cribb2025}, against the backdrop of a chronic productivity and business-investment gap \citep{williams2025}.

\subsection*{Gaps and open questions}

The synthesis leaves several gaps. Exposure measures are reported as point estimates whose dependence on the scoring model and prompt is acknowledged but rarely quantified, and none propagates within-model scoring uncertainty into the aggregates built on it. Sub-national exposure for the UK is estimated only at coarse geographies and has not been posed as a compositional question for a devolved fiscal authority, so the Scotland--rUK differential, and which part of it survives the measurement fragility above, is absent from the literature. And no instrument for AI exposure yet matches the commuting-zone design of \citet{acemoglu2020}, so credible sub-national causal estimates are unavailable, and what remains tractable in the near term is the measurement of the exposure differential and the propagation of its uncertainty into forecast implications. These gaps motivate three questions:

\begin{enumerate}
\item Which exposure estimands are robust to cross-model divergence? Specifically, do regional differentials survive the model-specific bias that destabilises absolute levels, and can that robustness be established analytically rather than asserted?
\item Does occupational composition insulate or expose a public-sector-intensive economy such as Scotland, once exposure is decomposed on the substitution--complementarity margin rather than read off a single exposure average?
\item How should the measurement uncertainty in LLM-based exposure scores be carried through to the productivity, employment, and wage projections that fiscal bodies such as the SFC and OBR rely on, so that a regional path is reported as a calibrated interval rather than a point?
\end{enumerate}

The remainder of this dissertation takes up these three questions in turn for the Scotland--rUK case.